% ============================================================
% BEAMER — GROUP 2 | Assignment 2
% Advanced Survival Analysis: Colon Cancer Dataset
% Non-Parametric, Semi-Parametric & Extended Cox Models
% Style : CambridgeUS, AIMS green — consistent with SA1 & SA2
% ============================================================
\documentclass[aspectratio=169, 10pt]{beamer}
\usepackage{xcolor} 
\usepackage{caption}
\usetheme{CambridgeUS}
\usecolortheme{default}
\usefonttheme{professionalfonts}
\captionsetup[figure]{labelformat=simple, labelsep=colon} 
\captionsetup[table]{labelformat=simple, labelsep=colon} 
% ── Colour palette ──────────────────────────────────────────
\definecolor{primary}{RGB}{30, 80, 55}   % AIMS dark green
\definecolor{accent} {RGB}{69,123,157}   % steel blue
\definecolor{steelblue}{rgb}{0.274, 0.51, 0.706}

\usecolortheme[named=primary]{structure}
%\setbeamercolor{title}               {fg=white,  bg=primary}
%\setbeamercolor{frametitle}          {fg=white,  bg=primary}
%\setbeamercolor{block title}         {fg=white,  bg=primary}
%\setbeamercolor{block body}          {fg=black,  bg=gray!7}
%\setbeamercolor{alerted text}        {fg=red!70!black}
%\setbeamercolor{palette primary}     {bg=primary,      fg=white}
%\setbeamercolor{palette secondary}   {bg=primary!75,   fg=white}
%\setbeamercolor{palette tertiary}    {bg=primary!60,   fg=white}
%\setbeamercolor{author in head/foot} {bg=primary!80,   fg=white}
%\setbeamercolor{title in head/foot}  {bg=primary,      fg=white}
%\setbeamercolor{date in head/foot}   {bg=primary!60,   fg=white}

\setbeamertemplate{caption}[numbered]
% ── Footer ──────────────────────────────────────────────────
\setbeamertemplate{footline}{\small
  \leavevmode\hbox{
    \begin{beamercolorbox}[wd=.35\paperwidth,ht=2.5ex,dp=1ex,center]
      {author in head/foot}\scriptsize Group 2 - AIMS Rwanda
    \end{beamercolorbox}%
    \begin{beamercolorbox}[wd=.3\paperwidth,ht=2.5ex,dp=1ex,center]
      {title in head/foot}
      \href{https://github.com/tchfare/Survival-Analysis-Group2}
           {\textcolor{steelblue}\scriptsize GitHub Repo}
    \end{beamercolorbox}%
    \begin{beamercolorbox}[wd=.335\paperwidth,ht=2.5ex,dp=1ex,right]
      {date in head/foot}
      
      \insertframenumber/\inserttotalframenumber\hspace*{2ex}
    \end{beamercolorbox}}%
}

% ── Fonts & bullets ─────────────────────────────────────────
\setbeamerfont{title}     {series=\bfseries, size=\Large}
\setbeamerfont{frametitle}{series=\bfseries, size=\large}
\setbeamertemplate{navigation symbols}{}
\setbeamertemplate{itemize item}   {\color{primary}\textbullet}
\setbeamertemplate{itemize subitem}{\color{primary}$\circ$}
\setbeamertemplate{enumerate item} {\color{primary}\insertenumlabel.}

% ── Packages ────────────────────────────────────────────────
\usepackage{booktabs}
\usepackage{graphicx}
\usepackage{xcolor}
\usepackage{amsmath}
\usepackage{caption}
\usepackage{url}
\usepackage{hyperref}
\usepackage{multirow}
\hypersetup{colorlinks=false, urlcolor=black, linkcolor=black, citecolor=black}
\captionsetup{font=scriptsize, labelfont=bf, skip=1pt, labelformat=empty}

% ── Figure path — update to match your local outputs folder ─
\graphicspath{
  {/home/guest/Desktop/SA/Project_1/outputs2/figures/}
}

% ── Title block ─────────────────────────────────────────────
\title{\textbf{Survival Analysis}}
\subtitle{Assignment 2\\[3pt]
  Non-Parametric, Semi-Parametric \& Extended Cox Models}
\author[Group 2]{%
  \small
  Tchandikou Ouadja Fare, \quad
  Rabecca Kanini Kating'u, \quad
  Joshua Pius Opio, \\[3pt]
  Djadida Uwituze, \quad
  Josiane Kazanenda
}
\date{\scriptsize\today}
\institute{African Institute for Mathematical Sciences (AIMS Rwanda)}

% ============================================================
\begin{document}
% ============================================================

% ── TITLE ───────────────────────────────────────────────────
\begin{frame}
  \titlepage
  \vspace{-10pt}
%  \begin{center}
%    \small\color{primary}
%    %\url{https://github.com/tchfare/Survival-Analysis-Group2}
%  \end{center}
\end{frame}

% ── OUTLINE ─────────────────────────────────────────────────
\begin{frame}{Outline}
  \tableofcontents[hideallsubsections]
\end{frame}



% INTRODUCTION FRAME
\begin{frame}
	\begin{center}
		\colorbox{white}{\fbox{\parbox{0.8\linewidth}{\centering \Huge \texttt{INTRODUCTION}}}}
	\end{center}
\end{frame}









% ============================================================
\section{INTRODUCTION}
% ============================================================

\begin{frame}{Objectives}

			\begin{block}{General Objective}
		\large
				Extend the parametric analysis of colon cancer patients using
				non-parametric and semi-parametric methods to obtain an assumption-robust
				characterisation of survival patterns and covariate effects.
				
			\end{block}
			
			\vspace{6pt}
			\begin{block}{Specific Objectives}
				%\small
				\begin{enumerate}
					\item Estimate survival empirically (Kaplan-Meier)
					\item Characterise cumulative hazard (Nelson-Aalen)
					\item Test group differences with \alert{multiple
						weighted tests} (Log-Rank, Breslow,
					Tarone-Ware, Fleming-Harrington)
					\item Model covariate effects via Cox PH
					\item Diagnose, refine, and \alert{extend} the
					Cox model (log-log plots, time interactions)
					\item Synthesise all findings
				\end{enumerate}
			\end{block}

	\end{frame}
	
	
\begin{frame}{Research Questions}	
	\begin{enumerate}
		\large
		\item What does the Kaplan-Meier
					curve reveal about overall survival?\\
					
		\item What shape does the cumulative
					hazard take (Nelson-Aalen)?\\
					
					\item Are survival differences
					confirmed across all weighted tests?
					
					\item Which covariates significantly
					affect the hazard (Cox PH)?
					
					\item Do log-log plots confirm PH,
					and what does the spline model reveal?
					
					\item \alert{Does the Lev+5FU
						treatment effect change over time?}
						
					\item How do all approaches
					converge?
				\end{enumerate}	
\end{frame}




\section{METHODOLOGY}


\begin{frame}
	\begin{center}
		\colorbox{white}{\fbox{\parbox{0.8\linewidth}{\centering \Huge \texttt{METHODOLOGY}}}}
	\end{center}
\end{frame}



\begin{frame}{About the data}
	
	
	\textbf{Dataset} \medskip
	
	The \texttt{colon} dataset (North Central Cancer Treatment Group) contains
	929 patients from a randomised adjuvant chemotherapy trial.
	Key variables: treatment (\texttt{rx}: Obs, Lev, Lev+5FU), age, sex,
	nodes (positive lymph nodes), and tumour differentiation.
	
	\begin{center}
		
		\begin{table}[ht]  
			% Wrap the table inside the table environment
			\caption{Summary of Patient Data in the Colon Cancer Dataset}  % Table caption placed above the table
			\begin{center}
				\begin{tabular}{lrr}
					\toprule
					& \textbf{n} & \textbf{\%} \\
					\midrule
					Total patients & 929 & 100.0 \\
					Events (deaths) & 452 & 48.7  \\
					Censored        & 477 & 51.3  \\
					\bottomrule
				\end{tabular}
			\end{center}
			\vspace{5pt}  % Adds vertical space after the table
			\begin{flushleft}
				\centering
				\small
				$\blacksquare$ The dataset contains 929 patients, with 452 events (deaths) and 477 censored observations.
			\end{flushleft}
		\end{table}
	\end{center}
	
\end{frame}


% ── SLIDE: Data Preparation ─────────────────────────────────
\begin{frame}{Data Preparation}
	

	\vspace{6pt}
	
	\begin{block}{}
		\begin{enumerate}
			\large
			\item Restricted to \textbf{death events only} (\texttt{etype = 2}),
			one record per patient
			\item \textbf{Missing data:} \texttt{nodes and differ} had $<2\%$ missing values,
			imputed using Last Observation Carried Forward (LOCF)
			\item \textbf{Factor encoding:} \texttt{sex} (Female/Male),
			\texttt{rx} (Obs/Lev/Lev+5FU),
			\texttt{differ} (Well/Moderate/Poor)
			\item \textbf{Survival object:} constructed as
			\texttt{Surv(time, status)} in R
		\end{enumerate}
	\end{block}
	
\end{frame}





% ── SLIDE: Step 3 Methods ───────────────────────────────────
\begin{frame}{Hypothesis Testing}
	
	\begin{block}{Log-Rank \& Wilcoxon Tests}
		\large
		To formally test for survival differences between groups,
		two non-parametric tests were applied:
		\[
		\chi^2 = \frac{\left(\sum_j w_j(O_{1j} - E_{1j})\right)^2}
		{\sum_j w_j^2 V_j}
		\]
		where $O_{1j}$ and $E_{1j}$ are observed and expected deaths
		in group~1 at time $t_j$, and $V_j$ is the hypergeometric
		variance.
	\end{block}
	
	\vspace{6pt}
	
	\begin{columns}[T]
		\begin{column}{0.48\textwidth}
			\large
			$\blacksquare$ \textbf{Log-Rank} ($w_j = 1$): equal weight
			at all time points; optimal under proportional hazards;
			sensitive to late differences.
		\end{column}
		\begin{column}{0.48\textwidth}
			\large
			$\blacksquare$ \textbf{Wilcoxon/Breslow} ($w_j = n_j$):
			weights early events more heavily; more sensitive to
			early survival differences.
		\end{column}
	\end{columns}
	
\end{frame}


% ── SLIDE: Step 4 Methods ───────────────────────────────────
\begin{frame}{Cox Proportional Hazards Model}
	
	\begin{block}{Cox PH Model}
		\large
		The Cox proportional hazards model relates covariates to the
		hazard without specifying the baseline hazard $h_0(t)$:
		\[
		h(t \mid \mathbf{x}) = h_0(t)\,
		\exp\!\bigl(\beta_1 x_1 + \beta_2 x_2 + \cdots + \beta_p x_p\bigr)
		\]
		Parameters are estimated via \textbf{partial likelihood}.
		The \textbf{hazard ratio} for covariate $x_k$ is
		$\text{HR}_k = e^{\hat\beta_k}$, interpreted as the
		multiplicative change in hazard per unit increase in $x_k$.
	\end{block}
	
	\vspace{6pt}
	
	\begin{block}{Modelling Strategy}
		\begin{enumerate}
			\large
			\item \textbf{Univariable models:} one covariate at a time
			— age, sex, rx, nodes, differ
			\item \textbf{Multivariable model:} all covariates
			simultaneously, yielding adjusted hazard ratios
		\end{enumerate}
	\end{block}
	
\end{frame}


% ── SLIDE: Step 5 Methods ───────────────────────────────────
\begin{frame}{Diagnostics \& Model Refinement}
	
	\begin{block}{PH Assumption \& Residual Diagnostics}
		\large
		Three diagnostic tools were applied to the multivariable model:
	\end{block}
	
	\vspace{4pt}
	
	\begin{enumerate}
		\small
		\item \textbf{Schoenfeld residuals} (\texttt{cox.zph}): test
		$H_0$: scaled residuals are uncorrelated with time.
		A significant $p$-value indicates PH violation.
		
		\vspace{4pt}
		\item \textbf{Martingale residuals}: plotted against continuous
		covariates (age, nodes) to detect non-linearity in the
		log-hazard. A non-flat LOESS trend indicates a violation
		of the linearity assumption.
		
		\vspace{4pt}
		\item \textbf{dfbeta residuals}: measure the influence of each
		observation on each coefficient. Values outside $\pm 0.02$
		flag influential patients.
	\end{enumerate}
		\vspace{4pt}
		
	$\blacksquare$ \textbf{Refined model:} where non-linearity was detected
		(nodes), a \textbf{penalised spline}
		(\texttt{pspline(nodes, df=4)}) was incorporated to
		flexibly model the log-hazard relationship.
	
	
\end{frame}


% ── SLIDE: Step 6 / Software ────────────────────────────────
\begin{frame}{Synthesis \& Software}
	
	\begin{block}{Synthesis}
		\small
		Results from all methods were integrated to answer whether the
		non-parametric, semi-parametric, and parametric (GenGamma,
		Project~1) approaches converge in their characterisation of:
		\begin{itemize}
			\large
			\item[$\blacksquare$] The shape of the hazard over time
			\item[$\blacksquare$] The magnitude and significance of
			treatment effects
			\item[$\blacksquare$] The key predictors of survival
		\end{itemize}
	\end{block}
	
	\vspace{6pt}
	
	\begin{block}{Software}
		\footnotesize
		All analyses were performed in \textbf{R} using the following
		packages:
		
		\vspace{4pt}
		\begin{columns}[T]
			
			\begin{column}{0.48\textwidth}
				\begin{itemize}
					\footnotesize
					\item[$\blacksquare$] \texttt{survival} — KM, NA,
					Cox PH, diagnostics
					\item[$\blacksquare$] \texttt{survminer} — KM plots
					with risk tables
				\end{itemize}
			\end{column}
			\begin{column}{0.48\textwidth}
				\begin{itemize}
					\footnotesize
					\item[$\blacksquare$] \texttt{flexsurv} — GenGamma
					parametric fit (Project 1)
					\item[$\blacksquare$] \texttt{ggplot2} — Nelson-Aalen
					and synthesis figures
				\end{itemize}
			\end{column}
		\end{columns}
	\end{block}
	
\end{frame}

































\section{RESULTS AND INTERPRETATION}





\begin{frame}
	\begin{center}
		\colorbox{white}{\fbox{\parbox{0.8\linewidth}{\centering \Huge \texttt{RESULTS AND INTERPRETATION}}}}
	\end{center}
\end{frame}



% ============================================================
%\section{Kaplan-Meier Estimator}
% ============================================================

\begin{frame}{Kaplan-Meier Estimator: Overall Survival}

\begin{columns}[T]
  \begin{column}{0.56\textwidth}
  	\begin{figure}
  		\centering
  		\includegraphics[width=\linewidth]{/home/guest/Desktop/SA/p2/outputs_proj2/figures/Figure1a_KM_Overall.png}
  		\caption{Overall Kaplan-Meier Survival Curve}
  		\label{fig:KM_Overall}
  	\end{figure}
  \end{column}
  
  
  \begin{column}{0.41\textwidth}
    \medskip
%    \textbf{Survival estimates at key time points:}
%
%    \smallskip
%    \begin{tabular}{@{}rrrr@{}}
%     
%      \toprule
%      \textbf{Days} & \textbf{S(t)} & \textbf{Lower} & \textbf{Upper}\\
%      \midrule
%       365 & 0.863 & 0.839 & 0.888 \\
%       730 & 0.745 & 0.716 & 0.775 \\
%      1095 & 0.663 & 0.632 & 0.695 \\
%      1825 & 0.577 & 0.545 & 0.612 \\
%      2555 & 0.481 & 0.447 & 0.518 \\
%      \bottomrule
%    \end{tabular}
%
%    \bigskip
 
    \smallskip
    \begin{table}[ht]  % Wrap the table inside the table environment
    	\caption{Survival Quantiles for Colon Cancer Patients}  % Table caption placed above the table
    	\begin{center}
    		\small
    		\begin{tabular}{@{}lrl@{}}
    			\toprule
    			\textbf{Quantile} & \textbf{Days} & \textbf{95\% CI} \\
    			\midrule
    			25th  & 1235 & [1106, 1410] \\
    			Median & 2418 & [2152, ---]  \\
    			75th   & ---  & ---           \\
    			\bottomrule
    		\end{tabular}
    	\end{center}
    	\vspace{5pt}  % Adds vertical space after the table
    \end{table}
  
  
  
    \bigskip
    \small
    
    $\blacksquare$ 86\% survive at 1~year.\\
    \vspace{0.3cm}
     $\blacksquare$ Median survival $\approx$~2\,418 days ($\approx$6.6~years),
    consistent with the GenGamma mean of 2\,130 days
    from Project~1.
  \end{column}
\end{columns}

\end{frame}

% ── KM by treatment & sex ────────────────────────────────────
\begin{frame}{Kaplan-Meier Estimator: Stratified by Group}

\begin{columns}[T]
  \begin{column}{0.49\textwidth}
  	\begin{figure}
    \includegraphics[width=0.9\linewidth]{/home/guest/Desktop/SA/p2/outputs_proj2/figures/Figure1b_KM_Treatment.png}
    \caption{Kaplan-Meier Curve Stratified by treatment}
    \label{fig:KM_Treatment}
	\end{figure}	
\vspace{-0.45cm}
    \small
    $\blacksquare$ Lev+5FU separates clearly from Obs and Lev
    after $\sim$500~days, indicating a sustained
    long-term benefit.
  \end{column}
  \begin{column}{0.52\textwidth}
  	\begin{figure}
    \includegraphics[width=0.9\linewidth]{/home/guest/Desktop/SA/p2/outputs_proj2/figures/Figure1c_KM_Sex.png}
    \caption{Kaplan-Meier Curve Stratified by Sex}
    \label{fig:KM_sex}
    \end{figure}
    	
\vspace{-0.45cm}
    \small
    $\blacksquare$ The female and male curves overlap throughout
    follow-up, suggesting no meaningful sex effect
    on survival.
  \end{column}
\end{columns}

\end{frame}

% ============================================================
%\section{Nelson-Aalen Estimator}
% ============================================================

\begin{frame}{Nelson-Aalen Estimator: Cumulative Hazard \& Hazard Rate}

\begin{columns}[T]
  \begin{column}{0.49\textwidth}
  	
  	\begin{figure}
    \includegraphics[width=\linewidth]{/home/guest/Desktop/SA/p2/outputs_proj2/figures/Figure2a_NelsonAalen_CumHaz.png}
    \caption{Nelson Aalen Cumulative Hazard}
    \label{fig:NA_Ht}
    \end{figure}
    \vspace{-0.45cm}
    
    \small
    $\blacksquare$ $H(t)$ is \textbf{concave}, it rises quickly
    at first then flattens, indicating decelerating
    cumulative risk.
  \end{column}
  \begin{column}{0.49\textwidth}
  	
  	\begin{figure}
    \includegraphics[width=\linewidth]{/home/guest/Desktop/SA/p2/outputs_proj2/figures/Figure2b_NelsonAalen_Hazard.png}
    \caption{Nelson Aalen (Non-Parametric) Hazard}
    \label{fig:NA_H}
    \end{figure}
    
    \vspace{-0.45cm}
    \scriptsize
    $\blacksquare$ The smoothed hazard $h(t)$ is \textbf{non-monotone}:
    it rises to a peak near day~300 then gradually
    declines. This matches the Generalised Gamma
    model from the 1st presentation ($Q = -0.481 < 0$),
    confirming a non-monotone hazard shape.
  \end{column}
\end{columns}

\end{frame}

% ============================================================
%\section{Hypothesis Testing}
% ============================================================

\begin{frame}{Log-Rank \& Wilcoxon Tests -- Hypothesis Testing}

\begin{columns}[T]
  \begin{column}{0.54\textwidth}
    \begin{table}[ht]
    	\centering
    	% Wrap the table inside the table environment
    	\caption{Log-Rank Vs Wilcoxon Tests for Treatment and Sex.} % Table caption placed above the table
    	\begin{center}
    		\small
    		\begin{tabular}{@{}llrrl@{}}
    			\toprule
    			\textbf{Comparison} & \textbf{Test} &
    			$\boldsymbol{\chi^2}$ & \textbf{p} &
    			\textbf{Decision} \\
    			\midrule
    			Treatment (rx) & Log-Rank  & 16.79 & $<$0.05 & Reject $H_0$ \\
    			Treatment (rx) & Wilcoxon  & 14.45 & $<$0.05 & Reject $H_0$ \\
    			Sex            & Log-Rank  &  0.40 & $>$0.05 & Fail to reject \\
    			Sex            & Wilcoxon  &  0.38 & $>$0.05 & Fail to reject \\
    			\bottomrule
    		\end{tabular}
    	\end{center}
    \end{table}
    
    

    \bigskip
    \textbf{Log-Rank vs Wilcoxon}

    \smallskip
    The Log-Rank test assigns equal weight to all event times
    and is optimal under proportional hazards. The Wilcoxon
    (Breslow) test weights early events more heavily.
    Both tests agree here, indicating the treatment effect
    is consistent across early and late follow-up.
  \end{column}

  \begin{column}{0.43\textwidth}
    \textbf{Key findings}
    \medskip
    \begin{itemize}
      \item [$\blacksquare$] \textbf{Treatment (rx):} Both tests reject $H_0$.
            Survival differs significantly across treatment
            groups. Lev+5FU has the best prognosis.
      \item [$\blacksquare$] \textbf{Sex:} Both tests fail to reject $H_0$.
            No significant survival difference between
            male and female patients.
      \item [$\blacksquare$] Agreement between both tests suggests the
            treatment effect is stable over time, not
            concentrated at early or late failure times.
    \end{itemize}
  \end{column}
\end{columns}

\end{frame}

% ============================================================
%\section{Cox Proportional Hazards Model}
% ============================================================

\begin{frame}{Cox Proportional Hazards: Univariable Models}

For each covariate, a univariable Cox model was fitted using
$h(t \mid x) = h_0(t)\,e^{\beta x}$.
The table reports the estimated hazard ratio (HR), 95\% CI,
and significance.

\bigskip
\begin{center}
\small
\begin{table}[ht]  % Wrap the table inside the table environment
	\caption{Univariable Cox Model with Hazard Ratios}  % Table caption placed above the table
	\begin{center}
		\small
		
		
		   \begin{tabular}{@{}llrrrrl@{}}
			\toprule
			\textbf{Covariate} & \textbf{Level} & \textbf{HR}
			& \textbf{Lower} & \textbf{Upper}
			& \textbf{p-value} & \textbf{Sig} \\
			\midrule
			Age (per year)   & ---      & 1.005 & 0.998 & 1.012 & 0.1490 & \\
			Sex              & Male     & 0.913 & 0.754 & 1.106 & 0.3620 & \\
			Treatment        & Lev      & 0.978 & 0.796 & 1.202 & 0.8290 & \\
			Treatment        & Lev+5FU  & 0.656 & 0.530 & 0.813 & \alert{$<$0.001} & \alert{***} \\
			Nodes (per node) & ---      & 1.079 & 1.066 & 1.093 & \alert{$<$0.001} & \alert{***} \\
			Differentiation  & Moderate & 1.261 & 0.893 & 1.781 & 0.1840 & \\
			Differentiation  & Poor     & 1.638 & 1.124 & 2.387 & \alert{0.0098} & \alert{**} \\
			\bottomrule
			%\multicolumn{7}{l}{\scriptsize *** $p<0.001$;\quad ** $p<0.01$;\quad * $p<0.05$}
		\end{tabular}
		
				
%		\begin{tabular}{@{}llrrrl@{}}
%			\toprule
%			\textbf{Covariate} & \textbf{Level} & \textbf{HR} & \textbf{Lower} & \textbf{Upper} & \textbf{Sig} \\
%			\midrule
%			Age (per year)   & ---          & 1.005 & 0.998 & 1.012 & \\
%			Sex              & Male         & 0.913 & 0.754 & 1.106 & \\
%			Treatment        & Lev          & 0.978 & 0.796 & 1.202 & \\
%			Treatment        & Lev+5FU      & 0.656 & 0.530 & 0.813 & *** \\
%			Nodes (per node) & ---          & 1.079 & 1.066 & 1.093 & *** \\
%			Differentiation  & Moderate     & 1.261 & 0.893 & 1.781 & \\
%			Differentiation  & Poor         & 1.638 & 1.124 & 2.387 & ** \\
%			\bottomrule
%		\end{tabular}
	\end{center}
	\vspace{5pt}  % Adds vertical space after the table
	\begin{flushleft}
		\centering
		\scriptsize
		\vspace{-0.2cm} 
		*** $p<0.001$;\quad ** $p<0.01$;\quad * $p<0.05$  % Footnote for significance levels
	\end{flushleft}
\end{table}
\end{center}
\vspace{-0.25cm}
\small
$\blacksquare$ Lev+5FU, nodes, and poor tumour differentiation are
significant univariable predictors. Age, sex, and Lev alone
are not significant.

\end{frame}

% ── Multivariable Cox ────────────────────────────────────────
\begin{frame}{Cox Proportional Hazards: Multivariable Model}

The multivariable model includes: age, sex, rx, nodes,
and differentiation simultaneously.

\bigskip
\begin{columns}[T]
  
  	\begin{column}{0.50\textwidth}
  		\begin{center}
  			\small
  			\begin{table}[ht]  % Wrap the table inside the table environment
  				\caption{Multivariable Cox Model with Hazard Ratios and Confidence Intervals}  % Table caption placed above the table
  				\begin{tabular}{@{}lrrrl@{}}
  					     \toprule
  					     \textbf{Covariate} & \textbf{HR}
  					       & \textbf{95\% CI} & \textbf{p-value} & \textbf{Sig} \\
  					     \midrule
  					     Age              & 1.004 & [0.996, 1.011] & 0.3250 & \\
  					     Sex (Male)       & 0.885 & [0.729, 1.074] & 0.2190 & \\
  					     Lev              & 0.973 & [0.793, 1.195] & 0.7970 & \\
  					     \alert{Lev+5FU}  & \alert{0.638} & \alert{[0.516, 0.789]} & \alert{$<$0.001} & \alert{***} \\
  					     \alert{Nodes}    & \alert{1.077} & \alert{[1.063, 1.091]} & \alert{$<$0.001} & \alert{***} \\
  					     Diff.\ Moderate  & 1.214 & [0.859, 1.716] & 0.2700 & \\
  					     \alert{Diff.\ Poor} & \alert{1.592} & \alert{[1.090, 2.325]} & \alert{0.0156} & \alert{**} \\
  					     \bottomrule
  					     \multicolumn{5}{l}{\scriptsize *** $p<0.001$;\quad ** $p<0.01$;\quad * $p<0.05$}
  					   \end{tabular}
  			\end{table}
  		\end{center}
  	\end{column}
  

  \begin{column}{0.47\textwidth}
    \textbf{Interpretation}
    \medskip
    \begin{itemize}
      \item [$\blacksquare$] \textbf{Lev+5FU:} HR = 0.638, a
            \alert{36\% reduction in hazard}
            vs.\ observation, adjusting for all other
            covariates.
      \item [$\blacksquare$] \textbf{Nodes:} Each additional positive
            lymph node raises hazard by 7.7\%.
      \item [$\blacksquare$] \textbf{Poor differentiation:} 59\% higher
            hazard vs.\ well-differentiated tumours.
      \item [$\blacksquare$] Age and sex remain non-significant after
            adjustment, consistent with univariable
            results.
    \end{itemize}
  \end{column}
\end{columns}

\end{frame}

% ============================================================
%\section{Diagnostics \& Refinement}
% ============================================================

\begin{frame}{PH Assumption: Schoenfeld Residuals (Diagnostics \& Refinement)}

\begin{columns}[T]
  \begin{column}{0.56\textwidth}
  	\begin{figure}
  	\includegraphics[width=\linewidth]{/home/guest/Desktop/SA/p2/outputs_proj2/figures/Figure3_Schoenfeld_Residuals.png}
    \caption{Schoenfeld Residuals Plots}
    \label{fig:SR}
	\end{figure}
	
  \end{column}
  \begin{column}{0.41\textwidth}
    \medskip
    \footnotesize
    The Schoenfeld residual test formally evaluates
    $H_0$: the scaled residuals are uncorrelated with
    time, i.e.\ the covariate effect is constant.

    \bigskip
    \begin{tabular}{@{}lrl@{}}
      \toprule
      \textbf{Covariate} & \textbf{p} & \textbf{PH?}\\
      \midrule
      Age    & $>$0.05 & \checkmark \\
      Sex    & $>$0.05 & \checkmark \\
      rx     & $>$0.05 & \checkmark \\
      Nodes  & $>$0.05 & \checkmark \\
      Differ & $>$0.05 & \checkmark \\
      Global & $>$0.05 & \checkmark \\
      \bottomrule
    \end{tabular}

    \bigskip
    \small
    $\blacksquare$ All $p > 0.05$. The PH assumption holds globally
    and for each covariate. No stratification or
    time-varying coefficients are required.
  \end{column}
\end{columns}

\end{frame}

% ── Linearity & Influence ────────────────────────────────────
\begin{frame}{Linearity Check \& Influential Observations (Diagnostics \& Refinement)}

\begin{columns}[T]
  \begin{column}{0.57\textwidth}
  	
  	\begin{figure}
    \includegraphics[width=\linewidth]{/home/guest/Desktop/SA/p2/outputs_proj2/figures/Figure4_Martingale_Residuals.png}
    \caption{Martingale Residuals against Age and Nodes}
    \label{fig:mr}
    \end{figure}
    
	\vspace{-0.4cm}
	
    \footnotesize
    $\blacksquare$ Martingale residuals plotted against continuous
    covariates.\\
    
    The LOESS (Locally Estimated Scatterplot Smoothing) trend for \texttt{nodes}
    is non-linear: a steep initial rise that flattens
    at higher node counts. This suggests the linearity
    assumption in the log-hazard is violated for nodes.
  \end{column}
  
  
  
  \begin{column}{0.40\textwidth}
  	\begin{figure}
    \includegraphics[width=\linewidth]{/home/guest/Desktop/SA/p2/outputs_proj2/figures/Figure5_Influential_Obs.png}
    \caption{Influential Observations}
    \label{fig:IO}
    \end{figure}
    
    \vspace{-0.2cm}
    
    \small
    \textbf{-- dfbeta} values for \texttt{nodes}.\\
    $\blacksquare$ All observations fall within $\pm 0.02$,
    indicating no single patient is distorting the results. The results are stable.
  \end{column}
  
\end{columns}

\end{frame}

% ── Refined Cox ──────────────────────────────────────────────
\begin{frame}{Refined Cox Model: Penalised Spline for Nodes}

\begin{columns}[T]
  \begin{column}{0.48\textwidth}
  	\begin{figure}
    \includegraphics[width=0.95\linewidth]{/home/guest/Desktop/SA/p2/outputs_proj2/figures/Figure6_Spline_Nodes.png}
    \caption{Non-Linear Effect of Nodes (Penalised Spline)}
    \label{PS}
    \end{figure}

	\vspace{-0.5cm}
    \footnotesize
    $\blacksquare$ The penalised spline (\texttt{pspline(nodes, df=4)}) 
    reveals a \textbf{non-linear} log-hazard relationship: risk rises 
    steeply from 0--5 positive nodes, then continues to increase 
    more gradually. The curve is truncated at the 95th percentile 
    to avoid unreliable estimates in the sparse tail.
    
  \end{column}
  
  \begin{column}{0.47\textwidth}
    \medskip
    \small
    \begin{table}[ht]  % Wrap the table inside the table environment
    	\caption{Refined Model}  % Table caption placed above the table
    	\begin{center}
    		\footnotesize
    		\begin{tabular}{@{}lrrrl@{}}
    			     \toprule
    			     \textbf{Covariate} & \textbf{HR}
    			       & \textbf{95\% CI} & \textbf{p-value} & \textbf{Sig} \\
    			     \midrule
    			     Age              & 1.004 & [0.997, 1.012] & 0.3190 & \\
    			     Sex (Male)       & 0.880 & [0.725, 1.069] & 0.2060 & \\
    			     Lev              & 0.970 & [0.790, 1.190] & 0.7650 & \\
    			     \alert{Lev+5FU}  & \alert{0.634} & \alert{[0.512, 0.784]} & \alert{$<$0.001} & \alert{***} \\
    			     pspline(nodes)   & \multicolumn{2}{l}{non-linear} & \alert{$<$0.001} & \alert{***} \\
    			     Diff.\ Moderate  & 1.212 & [0.858, 1.713] & 0.2750 & \\
    			     \alert{Diff.\ Poor} & \alert{1.584} & \alert{[1.085, 2.314]} & \alert{0.0170} & \alert{**} \\
    			     \bottomrule
    			   \end{tabular}
    	\end{center}
    	\vspace{5pt}  % Adds vertical space after the table
    \end{table}

    \small
    $\blacksquare$ The main conclusions are unchanged.
    The spline improves model fit without altering
    the clinical interpretation of Lev+5FU,
    nodes, or poor differentiation.
  \end{column}
\end{columns}

\end{frame}

% ============================================================
%\section{Synthesis}
% ============================================================

\begin{frame}{KM vs.\ Generalised Gamma (Synthesis)}

\begin{columns}[T]
  \begin{column}{0.58\textwidth}
  	
  	\begin{figure}
    \includegraphics[width=0.97\linewidth]{/home/guest/Desktop/SA/p2/outputs_proj2/figures/Figure10_KM_vs_Parametric.png}
    \caption{KM Vs Parametric (Generalised Gamma) Survival Curves}
    \label{KMP}
    \end{figure}
    
    
    \vspace{-0.5cm}
    
    
    \footnotesize
    $\blacksquare$ KM (solid) and Generalised Gamma (dashed) survival
    curves are nearly identical up to $t \approx 2\,000$
    days.  They diverge slightly only in the sparse
    upper tail, confirming excellent parametric fit.
  \end{column}
  \begin{column}{0.39\textwidth}
  	\small
    \medskip
    \textbf{Convergence across methods}
    \medskip
    \begin{enumerate}
      \item KM median ($\approx$2\,418 days) is close
            to the GenGamma mean (2\,130 days)
      \item Nelson-Aalen hazard peak ($\sim$day~300)
            matches the GenGamma non-monotone shape
            ($Q = -0.481$) exactly
      \item Log-Rank confirms treatment effect;
            Cox quantifies it: Lev+5FU,
            HR = 0.638
      \item Nodes is the strongest continuous
            predictor (HR $\approx$ 1.077 per node)
      \item Poor differentiation raises hazard
            by $\sim$59\% vs.\ well-differentiated
    \end{enumerate}
  \end{column}
\end{columns}

\end{frame}









\section{CONCLUSION}





\begin{frame}
	\begin{center}
		\colorbox{white}{\fbox{\parbox{0.8\linewidth}{\centering \Huge \texttt{CONCLUSION}}}}
	\end{center}
\end{frame}


% ── Conclusions ──────────────────────────────────────────────
\begin{frame}{Conclusions}
\large
\begin{itemize}\setlength\itemsep{6pt}  % Use itemize for square bullets
	\item[$\blacksquare$] \texttt{Overall survival:}
	86\% survive 1~year; median $\approx$~2\,418 days.
	Lev+5FU patients show markedly better long-term
	survival; sex has no significant effect.
	
	\item[$\blacksquare$] \texttt{Hazard shape:}
	Non-monotone cumulative hazard. The instantaneous
	hazard peaks near day~300 then declines,
	consistent with the GenGamma model ($Q = -0.481$).
	
	\item[$\blacksquare$] \texttt{Group differences:}
	Treatment significantly affects survival
	($p < 0.05$, both tests); sex does not.
	Log-Rank and Wilcoxon agree, confirming the
	treatment effect is stable over time.
	
	\item[$\blacksquare$] \texttt{Key predictors:}
	\texttt{Lev+5FU} (HR = 0.638, 36\% hazard
	reduction), \texttt{nodes} (HR = 1.077 per node,
	non-linear), and \texttt{poor differentiation}
	(HR = 1.592).
	
	\item[$\blacksquare$] \texttt{Synthesis:}
	Non-parametric (KM/NA), parametric (GenGamma),
	and semi-parametric (Cox) methods consistently
	identify a high early-risk phase, declining hazard,
	and a strong protective benefit from Lev+5FU.
	
\end{itemize}


\end{frame}

% ── Clinical recommendations ─────────────────────────────────
\begin{frame}{Clinical Implications}
\large
\begin{enumerate}\setlength\itemsep{8pt}

  \item \textbf{Prioritise Lev+5FU:} The 36\% hazard
        reduction is robust to covariate adjustment and
        confirmed by both non-parametric and
        semi-parametric methods.

  \item \textbf{Intensify post-operative monitoring in
        the first 300 days:} This is the period of
        highest hazard. Closer surveillance could improve
        early detection of recurrence.

  \item \textbf{Node count as a risk stratifier:}
        The non-linear effect means that even moving from
        0 to 3 positive nodes substantially elevates risk.
        Node status should guide adjuvant therapy decisions.

  \item \textbf{Poor tumour differentiation:}
        A 59\% higher hazard identifies this group as
        high-priority for aggressive treatment protocols.

\end{enumerate}

\end{frame}

% ── References ───────────────────────────────────────────────
\begin{frame}{References}

\small
\begin{itemize}\setlength\itemsep{5pt}
  \item Cox, D.R. (1972).
        Regression models and life-tables.
        \textit{J. Royal Statistical Society B}, 34(2), 187--220.

  \item Jackson, C.H. (2016).
        \texttt{flexsurv}: A platform for parametric survival
        modeling in R.
        \textit{Journal of Statistical Software}, 70(8), 1--33.

  \item Laurie, J.A.\ et al.\ (1989).
        Surgical adjuvant therapy of large-bowel carcinoma.
        \textit{Journal of Clinical Oncology}, 7(10), 1447--1456.

  \item Therneau, T.M.\ \& Grambsch, P.M.\ (2000).
        \textit{Modeling Survival Data: Extending the Cox Model}.
        Springer.

  \item Rocke, D.M.\ (2025).
        \textit{Parametric Survival Models} (Lecture notes).
        University of California, Davis.
\end{itemize}

\end{frame}

% ── Thank you ────────────────────────────────────────────────
\begin{frame}
  \begin{center}
  	\large
    \vspace{18pt}
    {\Large\color{primary}\textbf{Thank You}}\\[8pt]
    \large for your attention\\[20pt]
    \normalsize\textit{Survival Analysis - Assignment 2}\\[4pt]
    Group 2 --- AIMS Rwanda, \\
    March 2026\\[10pt]
    \small\color{gray}
    \url{https://github.com/tchfare/Survival-Analysis-Group2}
  \end{center}
\end{frame}

\end{document}
